% 18b-lineage.tex: что Полярис наследует, что сочетает и что может быть его собственным.
% Добавлен 25.09.2026, после замечания внешнего читателя: внутренняя прослеживаемость отчёта
% тщательна, а его научная родословная едва видна; до этого раздела он не цитировал ничего.
\section{Родословная: что унаследовано, а что нет}
\label{sec:lineage}

До этого раздела отчёт не цитировал ничего, и читатель не мог отличить изобретение от синтеза.
Этот раздел их разделяет. Утверждает он меньше, чем читатель мог бы предположить: в Полярис нет
новой криптографии, нового протокола и новой системы доказательств. Каждый примитив ниже чужой, и
дальше названо, чей.

\subsection{Унаследовано}

\begin{itemize}
  \item \emph{Подписи.} МЛ-ДСА и СЛХ-ДСА суть стандарты НИСТ \cite{fips204,fips205}. Путь миграции
    и его след аудита есть инженерия Полярис; стойкость примитивов нет, и ничто в этом репозитории
    на неё не влияет.
  \item \emph{Обязательства Меркла и прозрачность.} Дерево эпохи и пакеты якорей суть деревья Меркла
    \cite{merkle1988}. Журналы прозрачности следуют журналу только на добавление из прозрачности
    сертификатов, с доказательствами включения и согласованности \cite{rfc6962}, а соподпись
    свидетеля над головой журнала следует идее коллективной соподписи свидетелей \cite{syta2016}.
  \item \emph{Анонимная принадлежность.} Доказательство принадлежности зафиксированному множеству без
    указания, какой именно элемент, с нуллификатором, отказывающим повторному использованию, есть
    конструкция, которую сделал практичной Zerocash \cite{zerocash2014}; схема построена на Plonky2
    \cite{plonky2}. Более широкая цель, удостоверения, которые держатель показывает, не оставляя
    следа, восходит к Чауму \cite{chaum1985} и к системам анонимных удостоверений \cite{camlys2001},
    которых Полярис не реализует: это не общая система анонимных удостоверений
    (раздел~\ref{sec:privacy}).
  \item \emph{Предъявление и кошельки.} Верификатор для кошельков реализует OpenID4VP и его профиль
    высокой достоверности \cite{oid4vp,haip} поверх СД-ЖВТ и СД-ЖВТ ПУ \cite{sdjwt,sdjwtvc} со
    списком статусов токенов \cite{statuslist}. Путь мобильного документа следует ИСО/МЭК 18013-5
    \cite{iso18013}, а словарь удостоверений модели данных W3C \cite{vcdm2}.
  \item \emph{Доверяющие стороны.} Идентификаторы субъекта, свои для каждой доверяющей стороны, суть
    попарные идентификаторы OpenID Connect \cite{oidc}. Краткоживущее подписанное утверждение о
    статусе есть образец ответа OCSP, который сервер прикрепляет \cite{rfc6960}. Вход оператора
    использует WebAuthn \cite{webauthn2}.
  \item \emph{Удостоверение личности при зачислении.} Силы свидетельств и вывод уровня достоверности
    следуют НИСТ SP 800-63A \cite{sp80063a}; Полярис записывает прочтение и говорит, где это
    прочтение.
  \item \emph{Коды принуждения.} Второй секрет, сигнализирующий о принуждении, давно применяется в
    охранной сигнализации и банковском деле; изобретён он не здесь. Полярис добавляет измерение в
    разделе~\ref{sec:limitations}: против какого противника он хуже, чем ничего.
  \item \emph{Проверка тестов.} Мутационный анализ принадлежит ДеМилло, Липтону и Сейворду
    \cite{demillo1978}, обзор дали Цзя и Харман \cite{jia2011}.
\end{itemize}

\subsection{Что может быть собственным Полярис}

Насколько известно автору, и в порядке описания, а не притязания на приоритет:
\begin{itemize}
  \item \emph{устройство}: гарантии уровня прав держит схема, а не приложение, и каждой сопоставлена
    машинная проверка со своим тестом обнаружения;
  \item \emph{метод обеспечения достоверности}: мутация самого слоя обеспечения (проверок, триггеров,
    процедур, верификаторов и свидетелей), чтобы найти зелёные проверки, доказывающие не то
    утверждение, предмет таблицы~16 и исправлений этого отчёта;
  \item \emph{рабочий договор}: прогресс считается лишь как выжившая внешняя зависимость, с чертой
    закрытия, а не как внутренний рост.
\end{itemize}
Ни одно из этого не есть результат, который читателю следует принять на слово отчёта. Каждое
описано так, чтобы его можно было сравнить с предшествующими работами, которых этот раздел ещё не
цитирует, и читателя, знающего такие работы, просят их назвать.

\subsection{Чего этот раздел не решает}

Он не обозревает системы идентификации, программы цифровой личности или литературу об анонимных
удостоверениях; он называет лишь то, из чего механизмы Полярис непосредственно происходят. Сравнение
с развёрнутыми национальными системами есть таблица в README репозитория, где сказано, что каждая
из них развёрнута, а Полярис нет.
